% !TEX root = ../main.tex

\chapter{绪论}

\section{研究背景与问题提出}

互联网金融信息的增长改变了股票分析的输入形态。投资者在判断个股时，除了价格、成交量和财务报表，也会参考新闻快讯、论坛讨论、机构研报和搜索引擎中的热点信息。这些文本并不直接等同于可交易信号，但它们会影响市场预期的形成和扩散。大语言模型近几年在文本理解和信息整合方面进展较快，金融领域也从传统文本分类、情绪识别，逐步转向金融专用模型和金融智能体等方向。BloombergGPT 说明了大规模金融语料训练专用模型的可行性，FinGPT 则强调开放数据管线与轻量微调在金融场景中的价值；相关综述同时指出，实时性、幻觉、可解释性和评估体系仍是金融 LLM 应用中的主要问题\citep{lee2024finllms,li2023llmfinance,nie2024llmfinancialapplications,yang2023fingpt,wu2023bloomberggpt}。

股票市场高度受预期影响。价格变化不仅反映企业基本面，也反映投资者对未来事件的理解、情绪和博弈结果。相比结构化行情数据，文本信息往往更早承载预期变化：论坛评论可以反映散户讨论热度和短期情绪，新闻传递突发事件与公开信息，研报则提供相对专业的分析框架和逻辑链条。因此，本文把这些材料纳入同一预测流程，重点观察它们在短期方向任务中是否能补充 K 线难以表达的事件和预期线索。已有研究表明，社交媒体情绪与市场波动之间存在可观测关系，大语言模型也已表现出从新闻文本中识别未来市场反应的潜力\citep{bollen2011twitter,lopezlira2023chatgptstock}。

已有工作覆盖了金融情绪分析、检索增强生成、社交媒体预测、金融大模型和强化学习交易等方向。本文关注的不足主要有三点：其一，很多方法围绕单一信息源展开，例如只使用新闻、社交媒体或财报文本，和真实投研中同时查看多类材料的流程仍有距离；其二，许多实验建立在静态离线数据集上，对预测时的实时信息接入讨论较少；其三，论坛和社交媒体文本噪声较高，手工情绪指标或复杂结构模型在工程实现上并不总是方便复用。金融 LLM 的幻觉、可追溯性和评估标准问题也仍然存在\citep{lewis2020rag,ye2024r2ag,lee2024finllms,li2023llmfinance,nie2024llmfinancialapplications,zhou2025expertopinion}。

基于上述背景，本文不只讨论离线数据集上的方向分类，而是围绕预测时的信息获取流程设计系统。框架不预先写死复杂的情绪共识公式，而是把多源文本整理、证据归纳和趋势判断交给大语言模型完成，并通过在线页面把这一流程串联起来。

\section{研究目标与研究问题}

本文希望把个股查询时常见的几类材料放进同一条预测流程中。用户输入股票代码后，系统能够获取近期 Comments、News、Reports、外部检索结果和 K 线数据，再由大语言模型阅读这些材料并给出后续走势判断。这里的重点不是单独证明某一种文本源最有效，而是考察这些材料进入同一模型输入后，能否比只看行情数据提供更多判断依据。

围绕这一目标，需要解决的问题主要包括：评论、新闻、研报和网页检索结果怎样整理，模型才不至于被长文本和噪声干扰；金融文本中的正面线索、风险提示和分歧意见如何进入最终判断；外部文本和领域监督微调是否会改变未来 1 日、3 日和 7 日方向预测结果；以及在监督微调之后，是否可以继续用奖励信号优化模型的输出稳定性。

\section{本文的主要工作与贡献}

本文首先完成了股票分析材料的获取与整理流程。在真实使用场景中，系统会在预测时抓取最近的论坛评论、相关新闻、机构研报，以及通过 Web API 返回的补充网页信息。当前工程实现中，Web API 模块主要通过 SerpAPI 接入搜索引擎结果，用于补充公告、财报、股价异动和最新消息等内容；不同来源的文本在进入模型前会经过清洗、截断、排序和摘要。

其次，本文设计了基于大语言模型的分支总结和最终预测机制。系统会分别总结 Comments、News、Report 和 Web 检索结果，再通过统一提示模板生成未来趋势预测。在线输出会说明支持上涨的因素、可能的下行风险以及仍需留意的分歧，而不是只给出一个孤立标签。本文没有预先手工定义情绪共识公式，而是让模型在阅读摘要后完成观点归纳、一致性比较和冲突识别。

同时，本文开展了离线可行性验证实验。考虑到历史时点的搜索结果难以完整回放，离线测试采用未来 1 日、3 日和 7 日涨跌方向作为评价目标，重点比较微调前后、是否加入新闻和研报摘要这两类因素对结果的影响。根据目前实验结果，LoRA SFT 使 7 日 K 线单周期任务准确率从 49.47\% 提升到 52.18\%；在多周期任务中加入新闻和研报摘要后，排除提取失败样本的平均逐项准确率从 50.39\% 提升到 56.13\%。这些数字仍不足以支撑交易应用，但能说明该流程值得继续验证。

具体来说，本文的工作包括三点：一是实现一个能够实时接入 Comments、News、Reports 和搜索结果的股票分析流程，使预测输入不再局限于静态行情样本；二是设计分支 Summary 与综合预测 Prompt，把评论、新闻、研报和网页检索结果压缩为模型可处理的输入，减少对手工情绪指标的依赖；三是通过微调前后、有无外部文本的对照实验，观察领域微调和文本信息加入后对方向预测的影响，并补充 GRPO 多周期方向训练入口。

\section{论文结构安排}

本文共分为七章。第 1 章为绪论，介绍研究背景、问题提出、研究目标、主要工作与论文结构。第 2 章介绍相关技术与理论基础，包括股票短期走势预测任务、多源金融信息特点、大语言模型、Prompt 方法、监督微调与强化学习对齐方法。第 3 章设计面向股票预测的多源实时信息增强框架，说明系统整体流程、多源信息获取、预处理筛选、统一组织与预测输出设计。第 4 章介绍模型训练与优化方法，包括监督微调数据构建、训练设计以及强化学习优化思路。第 5 章介绍系统实现与在线部署方案，包括数据获取、Web API 调用、缓存存储、模型调用和结果解析。第 6 章给出实验设计与结果分析，比较不同模型、不同信息输入和不同预测周期下的效果。第 7 章总结全文工作，分析当前不足并展望后续研究方向。
